\documentclass[11pt,a4paper]{article}
\usepackage[T1]{fontenc}
\usepackage[utf8]{inputenc}
\usepackage{lmodern,amsmath,amssymb}
\usepackage[margin=25mm]{geometry}
\usepackage{longtable,booktabs,array,calc}
\usepackage[hidelinks]{hyperref}
\hypersetup{pdftitle={Galileo, Newton refinement and the Planck-scale question},pdfauthor={navstokgap research working notes}}
\newcommand{\tightlist}{\setlength{\itemsep}{0pt}\setlength{\parskip}{0pt}}
\setlength{\emergencystretch}{3em}
\setcounter{secnumdepth}{0}
\begin{document}
\hypertarget{galileos-falling-parabola-newtons-refinement-and-the-planck-scale-question}{%
\section{Galileo's falling parabola, Newton's refinement and the
Planck-scale
question}\label{galileos-falling-parabola-newtons-refinement-and-the-planck-scale-question}}

For horizontal inertial speed \(v\) and downward force \(F\), the area
between the inertial horizontal line and the falling parabola obeys

\[\boxed{\frac{3F}{v}A_{\rm inertial,fall}=\tau\Delta E,
\qquad \Delta E=\frac{F^2\tau^2}{2m}.}\]

This is the programme's geometric anchor. The physical question is what
prevents arbitrarily fine Newtonian refinement from retaining
distinguishable classical alternatives as this action-scaled area
shrinks. Planck's constant sets quantum phase resolution; a hard
prohibition on every smaller area or deterministic energy--time product
requires an additional argument.

N01, 2026-09-14. The geometric coefficients reuse C001/C027 and the
retained Principia derivation. Sections 3--4 add an exploratory closed
comparison of inertial and piecewise falling motion, with written review
and no ledger promotion. No novelty is asserted.

\hypertarget{which-newton-area-is-being-compared}{%
\subsection{1. Which Newton area is being
compared}\label{which-newton-area-is-being-compared}}

Take \(m,F,v,\tau>0\) and downward coordinate \(y\). The initial
horizontal inertial trajectory and the falling trajectory are

\[q_I(t)=(vt,0),\qquad
q_F(t)=\left(vt,\frac{Ft^2}{2m}\right),\qquad 0\le t\le\tau.\]

Close their geometric region by the vertical segment at \(x=v\tau\). Its
area and the gained vertical kinetic energy are

\[A_{\rm inertial,fall}=\int_0^{v\tau}\frac{Fx^2}{2mv^2}\,dx
=\frac{vF\tau^3}{6m},\qquad
\Delta E=F\Delta y=\frac{F^2\tau^2}{2m}.\]

For gravity put \(F=mg\). Thus the geometrical area is proportional to
\(\tau\Delta E\), with the indispensable factor \(3F/v\). Its units are
area; \(\tau\Delta E\) has action units. This is the parabola and
limiting-area geometry of the projectile Scholium and Lemmas X--XI in
the retained \href{principia-constant-force-action.md}{Principia note}.
Kepler's equal-area law and sectors swept about a force centre are a
separate construction.

Three elementary regions occur in this diagram:

\begin{longtable}[]{@{}ll@{}}
\toprule\noalign{}
Region & Area \\
\midrule\noalign{}
\endhead
\bottomrule\noalign{}
\endlastfoot
Inertial horizontal line to falling parabola & \(vF\tau^3/(6m)\) \\
Matched-endpoint chord to falling parabola & \(vF\tau^3/(12m)\) \\
Inertial horizontal line to endpoint chord & \(vF\tau^3/(4m)\) \\
\end{longtable}

The first area is twice the second; the third is their sum. The
programme's anchor is the first. The second is useful for comparing
actions with the same endpoints, but must not silently replace it.

For \(\mathcal L=m|\dot q|^2/2+Fy\), the existing C001 calculation gives

\[D_{\rm chord,fall}=S[q_{\rm chord}]-S[q_F]
=\frac{F^2\tau^3}{24m}
=\frac{\tau\Delta E}{12}
=\frac{F}{4v}A_{\rm inertial,fall}.\]

Both paths in this last comparison have the same endpoint positions and
times, so adding a total time derivative to the Lagrangian cancels from
their action difference. The inertial horizontal path instead ends at a
different height. Its action comparison needs endpoint/control phases
before it is a measurable relative phase. Section 3 supplies an explicit
reunion rather than hiding that requirement in the area-to-action
coefficient.

\hypertarget{the-precise-unresolved-planck-step}{%
\subsection{2. The precise unresolved Planck
step}\label{the-precise-unresolved-planck-step}}

At fixed \(m,F,v\), these areas and action products scale as \(\tau^3\).
Inserting arbitrarily small mathematical intervals makes them approach
zero. The research objective is a \textbf{physical obstruction to
continuing that refinement as distinguishable classical motion}, or a
stronger admissibility principle if one can be derived. Repeating the
classical zero limit does not complete this objective.

The deterministic gain \(\Delta E\) above is not a quantum state's
energy standard deviation. Total mechanical energy is conserved during
free fall; the kinetic gain is compensated by potential loss. Therefore
a statement about energy uncertainty and an observable evolution time
cannot be applied by simply renaming its two variables \(\Delta E\) and
\(\tau\).

With the quantum rule \(e^{iS/\hbar}\), small action differences give
nearby phases. Feynman's
\href{https://www.feynmanlectures.caltech.edu/II_19.html}{least-action
discussion} explicitly sums contributions from neighbouring paths,
including paths with action differences within \(\hbar\). It does not
impose a smallest allowed path-action difference. This supports the
breakdown of a uniquely resolved classical alternative, rather than an
automatic cutoff on the continuum of paths. A theorem forbidding
physical area shrinkage is a stronger target.

\hypertarget{close-the-inertial-and-falling-paths-with-specified-forces}{%
\subsection{3. Close the inertial and falling paths with specified
forces}\label{close-the-inertial-and-falling-paths-with-specified-forces}}

Consider a horizontal reference arm with zero transverse force and a
signal arm with the following controlled transverse force. Set
\(a=F/m\):

\[f(t)=\begin{cases}
 F,&0<t<\tau,\\
-F,&\tau<t<3\tau,\\
 F,&3\tau<t<4\tau.
\end{cases}\]

Both arms start at \(y=0\) with \(\dot y=0\). The signal arm's first
interval is exactly the falling parabola of Section 1. The remaining
intervals return it to the reference arm with the same position and
velocity. This is an ideal branch-controlled force experiment; a gravity
realization needs compensating forces on the reference and return
stages. It is not unassisted free fall throughout the whole comparison.

Let \(P(t)=\int_0^t f(s)\,ds\) and \(Y(t)=m^{-1}\int_0^t P(s)\,ds\).
Then \(P(4\tau)=Y(4\tau)=0\). Explicitly \(P/F\) is \(t\) on the first
interval, \(2\tau-t\) on the second, and \(t-4\tau\) on the third.
Consequently

\[\int_0^{4\tau}P(t)^2dt=\frac43F^2\tau^3,\qquad
\int_0^{4\tau}Y(t)dt=\frac{2F\tau^3}{m}.\]

\(Y\) stays nonnegative and reaches \(F\tau^2/m\) at \(t=2\tau\). At
common horizontal speed \(v\) the area enclosed between the signal
trajectory and its inertial baseline is

\[A_{\rm loop}=v\int_0^{4\tau}Y(t)dt
=\frac{2vF\tau^3}{m}=12A_{\rm inertial,fall}.\]

The transverse signal Lagrangian is \(m\dot Y^2/2+f(t)Y\); the reference
transverse Lagrangian vanishes. Since \(m\ddot Y=f\) and the endpoint
boundary term \(mY\dot Y\) vanishes,

\[\Delta S_{\rm loop}
=-\frac1{2m}\int_0^{4\tau}P(t)^2dt
=-\frac{2F^2\tau^3}{3m}
=-\frac{F}{3v}A_{\rm loop}
=-\frac43\tau\Delta E.\]

The sign and coefficient differ from the off-shell chord comparison,
because the return forces are part of this experiment. The exact
relation connects the original falling-parabola area to a closed,
dynamically specified phase comparison:
\(|\Delta S_{\rm loop}|=4F A_{\rm inertial,fall}/v\).

\hypertarget{the-same-closed-comparison-in-quantum-mechanics}{%
\subsection{4. The same closed comparison in quantum
mechanics}\label{the-same-closed-comparison-in-quantum-mechanics}}

Assume canonical quantum kinematics \([y,p]=i\hbar\), a coherent two-arm
label, and conditional Hamiltonians

\[H_I=\frac{p^2}{2m},\qquad H_F(t)=\frac{p^2}{2m}-f(t)y.\]

These assumptions supply \(\hbar\) and the phase rule; this is a test of
the Planck-scale interpretation, not a derivation of quantum mechanics
from classical geometry. No additional branch-dependent scalar
potentials or uncontrolled arm-label phases are included in the model.
Horizontal evolution is common and cancels in the relative propagator.

In the interaction picture of \(H_I\), the evolution generator is
\(B(t)=if(t)(y+tp/m)/\hbar\). Its commutator is the scalar

\[[B(t),B(s)]=\frac{i f(t)f(s)(t-s)}{m\hbar}.\]

All higher nested commutators vanish. The first integrated term vanishes
because \(\int f=\int tf=0\). The second gives the exact relative
unitary

\[U_I(4\tau)^\dagger U_F(4\tau)
=\exp\left[-\frac{i}{2m\hbar}\int_0^{4\tau}P(t)^2dt\right]I
=e^{i\Delta S_{\rm loop}/\hbar}I.\]

For the sign, integrate
\(\int f(t)\int_0^t f(s)(t-s)\,ds\,dt=-\int P(t)^2dt\); \(P(4\tau)=0\)
eliminates the boundary term. The algebra can first be done on Schwartz
states; the resulting forced free-particle unitaries extend the identity
to the motional Hilbert space. Thus the two wavepackets reunite exactly
for arbitrary common initial motional preparation, not just their
classical centres. A coherent balanced arm label then carries relative
phase \(\phi=\Delta S_{\rm loop}/\hbar\).

Use the already reviewed C006 comparison between this phase state and
the zero-phase state, with equal hypothesis priors. For \(n\)
independent copies and any joint binary quantum measurement, the success
probability is

\[p_{\rm opt}=\frac12\left(1+\sqrt{1-\cos^{2n}(\phi/2)}\right).\]

Within the first lobe \(|\phi|\le\pi\), a required success probability
\(p_*>1/2\) gives the action threshold

\[d_{n,p_*}=2\hbar\arccos\left([4p_*(1-p_*)]^{1/(2n)}\right),\qquad
\tau\ge\left(\frac{3m d_{n,p_*}}{2F^2}\right)^{1/3}.\]

Equivalently, in this specified comparison the original single falling
lens must have \(A_{\rm inertial,fall}\ge v d_{n,p_*}/(4F)\) for that
success rate. At \(n=1,p_*=3/4\), \(d=\pi\hbar/3\). For fixed \(p_*<1\)
it tends to zero as \(n\to\infty\); perfect discrimination instead
requires \(d=\pi\hbar\) for every finite \(n\). Phase periodicity
forbids extending the inequality beyond the first lobe as a monotone
criterion.

This is a positive, protocol-dependent obstruction to resolving the
shrinking Galileo area at fixed resources. It is not an absolute
prohibition on smaller loops: their quantum states remain defined and
approach the zero-phase state. It also does not assert that this
interference experiment optimizes every possible measurement of the
force or trajectory.

\hypertarget{research-consequence-and-next-decision}{%
\subsection{5. Research consequence and next
decision}\label{research-consequence-and-next-decision}}

The double-check recovers the intended chain: \textbf{inertial line
versus falling parabola -\textgreater{} action-scaled area
-\textgreater{} Planck-sensitive distinguishability}. It identifies an
actual drift source in STATE's substitution of Kepler swept sectors and
in README's stale audit queue. These entry points are corrected. The
arc--chord area remains a secondary matched-endpoint comparison with a
factor of two; it must not replace the user's inertial--parabola anchor.

The closed-force construction supplies a quantum benchmark. The user's
intended innovation is stronger: \textbf{a mathematical necessity
argument} that excludes zero-action refinement without assuming quantum
mechanics, an uncertainty relation, a positive action unit or a fixed
measurement budget. The quantum calculation above therefore supports the
interpretation but does not discharge that objective.

N02 must construct and decisively test one independently justified
consistency principle for joint time/position refinement of the Galileo
comparison. State the cut state, composition and limiting variables.
Derive the zero-action contradiction, or give an exact countermodel
satisfying the proposed principle. C027--C029 and the full-state
classical cuts already show why unqualified appeals to refinement
consistency do not suffice; a new principle must do more than repeat
those requirements. This returns to I003's originating
logical-obstruction question with its existing corrections retained. No
mathematical contradiction follows merely from the fact that nature is
quantum.
\end{document}
